\scanpage{27}
\mspage{8}
On prend un autom. \inserted{$u$} du groupoïde $\widehat{\mathrm{II}}_J$ qui \inserted{commute à $G_J$ et qui} fixe les sommets, et qui comme plus haut laisse stables les sous-groupoïdes $\widehat{\mathrm{II}}_{J,i}$ $(i\in J)$, \deleted{\unsure{engendrés|les groupoïdes}} $\mathrm{II}_{J,i}$ engendré par les \inserted{2} sommets $R_i^\omega$ $(\omega\in\uline{\omega})$ et les deux flèches $a_i^\omega$ $(\omega\in\uline{\omega})$.
\deleted{\unsure{Alors}} L'endom. $u$ est connu quand on connaît les $u(a_i^\omega)$, qui sont de la forme
\[
(20)\qquad u(a_i^\omega)=a_i^\omega\varphi_i^\omega\bigl(\underbrace{l_0^\gamma}_{\lambda}\bigr)
\qquad\begin{array}{l}\gamma=\frac{p-1}{2}\in\widehat{\mathbb Z}\\p\in\widehat{\mathbb Z}^{*}\end{array}
\]
\margin{NB: $\lambda\in\widehat\Pi_{0,3}$, pas la peine d'exiger de prime abord que $\lambda$ de la forme $l_0^\gamma\ldots$}
\uline{et} les
\[
(21)\qquad u(c_i^\omega)=\varphi_{\omega'i}^{\omega'}(\eta)c_i^\omega=c_i^\omega\psi_i^\omega(\eta)
\qquad\eta\in\widehat\Pi_{0,3}\quad\text{indépendant de }i,\omega
\]
On aura alors (comparer $(11)$)
\[
(22)\qquad u(b_i^\omega)=b_i^\omega\varphi_{\omega i}^\omega(\beta)\qquad\text{avec}
\]
\[
(23)\qquad\beta=\sigma_\infty(\eta)\eta^{-1}
\]
\margin{NB devient plus joli si $\eta$ est remplacé par son $\sigma_\infty\eta$, ou $\beta$ par $\beta^{-1}$}
Les relations à satisfaire \deleted{\unsure{sont}} alors déjà écrites en \inserted{$p$ (ou $\gamma$) et} $\beta$, \deleted{\unsure{la}} dans $(12)$ --- mais la première relation $\beta\sigma_\infty(\beta)=1$ conséquence de $(23)$, il reste la deuxième qui s'écrit donc
\margin{\uline{NB} Si $u$ remplacé par $u'=\tau(u)$, alors $\eta$ est remplacé par $\eta'=\tau_\infty(\eta)$, $\beta$ par $\tau_\infty(\beta)$, mais $\lambda$ par $\varepsilon_0^{-1}(\tau_\infty(\lambda^{-1}))$.}
\[
(24)\qquad\deleted{\unsure{(l_1^{-\gamma}\sigma_\infty(\eta)\eta^{-1}}}
\quad\alpha\rho^{-1}(\alpha)\rho^{-2}(\alpha)=1
\qquad\text{avec}\quad\alpha=\underbrace{l_1^{-\gamma}}_{\sigma_\infty(\lambda)}\underbrace{\sigma_\infty(\eta)\eta^{-1}}_\beta.
\]
On aura ici
\[
(25)\qquad
u(\widetilde a_i^\omega)=\psi_{\unsure{\omega i}}^{\unsure{\omega}}\bigl(\underbrace{\eta^{-1}l_0^\gamma\varepsilon_0(\eta)}_{\widetilde\alpha=\eta^{-1}\varepsilon_0(\lambda\eta)}\bigr)\widetilde a_i^\omega
=\widetilde a_i^\omega\psi_{\omega'i}^{\omega'}\bigl(\underbrace{\varepsilon_0^{-1}(\eta^{-1})l_0^\gamma\eta}_{\varepsilon_0^{-1}(\widetilde\alpha)}\bigr)
\]
\inserted{valable quel que soit $\lambda$}
\margin{$\widetilde\alpha=\eta^{-1}l_0^\gamma\varepsilon_0(\eta)$}
et
\[
(25\text{ bis})\quad\deleted{\unsure{\alpha=}}\quad\alpha=\rho\bigl(\eta\widetilde\alpha\rho^{-1}(\eta)^{-1}\bigr)
\]
\margin{valable pour $\lambda$ quelconque, pas \unsure{néc.} de la forme $l_0^\gamma$.}
\deleted{\unsure{car}} dans l'équation $(24)$ en $\alpha$ (invariante par substitutions $\alpha\mapsto\rho^{-1}(\alpha)$ (ou $\alpha\mapsto\rho\alpha=\rho^{-2}(\alpha)$) et $\alpha\mapsto x\alpha\rho^{-1}(x)^{-1}$,
\margin{26}
